\chapter{Respuesta en frecuencia de sistemas LIT TD}

\section{Respuesta en frecuencia (tiempo discreto)}

\subsection{Respuesta de un sistema LIT a una entrada peri\'{o}dica}

\begin{enumerate}

\item Si la entrada a un sistema de tiempo discreto es $x\left[ n\right] = z^n,$
entonces la salida es $y\left[ n\right] =H\left( z\right) z^n$ donde
\begin{equation}
  H\left( z\right) =\sum_{k=-\infty }^{\infty }h\left[ k\right] \, z^{-k},
\end{equation}
en la cual $h\left[ n\right] $ es la respuesta al impulso del sistema LIT.

\item En el caso general de $z$ complejo, $H\left( z\right) $ es la funci\'{o}n
de transferencia del sistema. 

\item En el caso $\left| z\right| =1\rightarrow z=e^{j \, \omega}$ y 
$x\left[ n\right] =z^n=e^{j \, \omega \, n},$ la respuesta en frecuencia del
sistema esta definida por 
\begin{equation}
H\left( e^{j \, \omega }\right) =\sum_{n=-\infty }^{\infty }h\left[ n\right]e^{-j \, \omega \, n}.
\end{equation}

\item La salida $y\left[ n\right] $ para una entrada 
\begin{equation}
x\left[ n\right] =\sum_{k=0}^{N-1}a_{k}e^{jk\left( \frac{2\pi }{N}\right) n},
\end{equation}
est\'{a} dada entonces por 
\begin{equation}
  y\left[ n\right] =
  \sum_{k=0}^{N-1}a_{k}H\left( e^{\frac{j2\pi k}{N}}\right)e^{jk\left( \frac{2\pi }{N}\right) n}.
\end{equation}

\end{enumerate}

\subsubsection{Ejemplo}

\begin{enumerate}

\item Sea un sistema con respuesta al impulso 
\begin{equation}
h\left[ n\right] =\alpha ^nu\left[ n\right] ,-1<\alpha <1,
\end{equation}
y una entrada 
\begin{equation}
x\left[ n\right] =
  cos\left( \frac{2 \, \pi \, n}{N}\right) =
  \frac{1}{2}\left( e^{j\frac{2 \, \pi \, n}{N}}+e^{-j\frac{2 \, \pi \, n}{N}}\right) .
\end{equation}

\item Calculamos la respuesta en frecuencia planteando la suma
\begin{equation}
  H\left( e^{j \, \omega }\right) =\sum_{n=-\infty }^{\infty }h\left[ n\right]e^{-j \, \omega \, n}.
\end{equation}
El sistema dado es causal porque la respuesta al impulso $\delta[n]$ es nula para $n < 0$. Por lo
tanto los l\'{\i}mites de la sumatoria cambian y resulta
\begin{equation}
H\left( e^{j \, \omega }\right) =
  \sum_{n=-\infty }^{\infty }\alpha ^nu\left[ n\right] e^{-j \, \omega \, n}=
  \sum_{n=0}^{\infty }\alpha ^ne^{-j \, \omega \, n}.
\end{equation}
Resolviendo la suma podemos escribir
\begin{equation}
H\left( e^{j \, \omega }\right) =
  \sum_{n=0}^{\infty }\alpha ^ne^{-j \, \omega}=
  \sum_{n=0}^{\infty }\left( \alpha e^{-j \, \omega }\right) ^n=
    \frac{1}{1-\alpha e^{-j \, \omega }}.
\end{equation}

\item La salida $y\left[ n\right] $ resulta ser 
\begin{equation}
y\left[ n\right] =
  \frac{1}{2} \left( 
     H\left( e^{j\frac{2 \, \pi \, n}{N}}\right) \, e^{j\frac{2 \, \pi \, n}{N}}+
     H\left( e^{-j\frac{2 \, \pi \, n}{N}}\right) \, e^{-j\frac{2 \, \pi \, n}{N}}\right) ,
\end{equation}
o sea 
\begin{equation}
y\left[ n\right] =
  \frac{1}{2}\left( 
     \frac{1}{1-\alpha \, e^{-j\frac{2 \, \pi \, n}{N}}} \, e^{j\frac{2 \, \pi \, n}{N}}+
     \frac{1}{1-\alpha \, e^{j\frac{2 \, \pi \, n}{N}}} \, e^{-j\frac{2 \, \pi \, n}{N}}
  \right) .
\end{equation}

\end{enumerate}

\subsection{Respuesta de un sistema LIT estable a entradas con transformada $Z$}

\begin{enumerate}

\item Supongamos un sistema LIT estable. Recordemos que en este caso, existe la transformada $Z$ de la
respuesta al impulso
\begin{equation}
  h[n] \begin{array}{c} Z \\ \leftrightarrow \end{array} H(z)
\end{equation}
y su regi\'{o}n de convergencia incluye el c\'{\i}rculo unitario. Esto garantiza adem\'{a}s la existencia
de la transformada de Fourier $H(e^{j \omega})$.

\item Sabemos que las transformadas $Z$ de la entrada
\begin{equation}
  x[n] \begin{array}{c} Z \\ \leftrightarrow \end{array} X(z)
\end{equation}
la salida
\begin{equation}
  y[n] \begin{array}{c} Z \\ \leftrightarrow \end{array} Y(z)
\end{equation}
y la respuesta al impulso 
de un sistema LIT est\'{a}n relacionadas por la expresi\'{o}n
\begin{equation}
   Y(z) = H(z) \, X(z).
\end{equation}
Calculando la antitransformada $y[n]$ obtenemos la solución de la ecuación de diferencias.

\end{enumerate}
\section{Representaci\'{o}n gr\'{a}fica de la respuesta en frecuencia de un sistema LIT}

La respuesta en frecuencia es una funci\'{o}n compleja y debe ser representada gr\'{a}ficamente en
\begin{enumerate}
        \item   Forma polar: magnitud y fase.
        
        \item   Forma cartesiana: parte real y parte imaginaria.
\end{enumerate}

Ambas representaciones permiten visualizar distintas propiedades de los sistemas LIT. Por tal motivo las
estudiaremos por separado.

\subsection{Definiciones}

\subsubsection{Ganancia o atenuaci\'{o}n}

La ganancia y la atenuaci\'{o}n de un filtro son conceptos complementarios, y sus valores est\'{a}n 
determinados por la magnitud de la respuesta en frecuencia. Dado que la salida de un filtro $y$ para
una entrada peri\'{o}dica $x$ queda determinada por 
\begin{equation}
        y[n] = H(e^{j \omega}) x[n],
\end{equation}
para tiempo discreto, el m\'{o}dulo de la funci\'{o}n compleja $\parallel H \parallel$ determina la 
variaci\'{o}n del tama\~{n}o de la salida respecto a la entrada. Hablamos de ganancia cuando deseamos 
saber cuanto queda ampliada la salida respecto a la entrada, y hablamos de atenuaci\'{o}n cuando deseamos 
saber cuanto queda disminuida la salida respecto a la entrada. Dado que por lo general
los filtros amplian determinadas frecuencias y disminuyen otras, hablar de ganancia o atenuaci\'{o}n es 
equivalente a elegir un punto de vista ($G = 1 - A$).

\subsubsection{Fase}

La fase de un filtro es la fase de la respuesta en frecuencia (funci\'{o}n compleja) $H(j \omega)$ en 
tiempo continuo o $H(e^{j \omega})$ para tiempo discreto.

\subsubsection{Fase lineal y no lineal}

Cuando la fase del filtro es una funci\'{o}n \textbf{lineal} de $\omega$, hay una interpretaci\'{o}n 
simple de su efecto en el dominio del tiempo. 
En el caso de tiempo discreto, tenemos sistemas con fase lineal cuando la pendiente de la fase con la 
frecuencia es un n\'{u}mero 
\textbf{entero}.
El sistema de tiempo discreto con respuesta en frecuencia
\begin{equation}
        H( e^{j \omega}) = e^{j \omega \, n_{0}},
\end{equation}
tiene una fase lineal 
\begin{equation}
        \angle H( e^{j \omega}) = -\omega \, n_{0},
\end{equation}
La salida de este sistema es igual a su entrada desplazada en el tiempo una cantidad \textbf{entera} de 
muestras $n_{0}$
\begin{equation}
        y[n] = x[n - n_{0}].
\end{equation}
Cuando $n_{0}$ no es un entero el comportamiento es m\'{a}s complejo y se explica m\'{a}s adelante.

Cuando la fase del filtro es una funci\'{o}n \textbf{no lineal} de $\omega$, cada componente exponencial 
compleja de la entrada (con diferentes frecuencias) son afectadas de forma tal que las fases relativas 
cambian y as\'{\i} obtenemos una salida que puede diferir en forma significativa respecto a la entrada.

\subsubsection{Retraso de grupo}

\section{Gr\'{a}ficos de Bode}

\subsection{Sistema de tiempo discreto con un polo}

Consideremos el sistema de tiempo discreto lineal, invariante en el tiempo y
causal cuya ecuaci\'{o}n de diferencias es
\begin{equation}
y\left[ n\right] -ay\left[ n-1\right] =x\left[ n\right] ,
\end{equation}
donde $\left| a\right| <1$.

\begin{enumerate}
\item  La funci\'{o}n de transferencia de este sistema es
\begin{equation}
H\left( z\right) =\frac{1}{1-az^{-1}}.
\end{equation}

\item  La respuesta en frecuencia de este sistema es
\begin{equation}
H\left( e^{j \, \omega }\right) =\frac{1}{1-ae^{-j \, \omega }}.
\end{equation}

\item  La respuesta al impulso es
\begin{equation}
h_{1}\left[ n\right] =a^{n}u\left[ n\right] .
\end{equation}

\item  La respuesta al escal\'{o}n es
\begin{equation}
h_{2}\left[ n\right] =h_{1}\left[ n\right]  u\left[ n\right] =\frac{%
1-a^{n+1}}{1-a}u\left[ n\right] .
\end{equation}

\item  La magnitud de $\left| a\right| <1$ determina la velocidad a la que
responde el sistema. 

\begin{enumerate}
\item  El sistema decae m\'{a}s r\'{a}pido para valores de $\left| a\right| $ cercanos a $0.$

\item  El sistema decae m\'{a}s lentamente para valores de $\left| a\right| $ cercanos a $1.$
\end{enumerate}

\item  Para valores de $-1<a<0,$ 

\begin{enumerate}
\item  la respuesta al impulso toma valores positivos y negativos alternados y decrecientes 
(oscilaciones amortiguadas). 

\item  la respuesta al escal\'{o}n tiene sobrepaso de su valor final y oscilaciones amortiguadas. 
\end{enumerate}

\item  El cuadrado de la magnitud de la respuesta en frecuencia es
\begin{equation}
H\left( e^{-j \, \omega }\right) H\left( e^{j \, \omega }\right) =
\frac{1}{1-ae^{j \, \omega }}\frac{1}{1-ae^{-j \, \omega }},
\end{equation}
que puede escribirse, usando la identidad de Euler, como
\begin{equation}
H\left( e^{-j \, \omega }\right) H\left( e^{j \, \omega }\right) =
\frac{1}{1-a\left( \cos \left( \omega \right) +j\sin \left( \omega \right) \right) }
\frac{1}{1-a\left( \cos \left( \omega \right) -j\sin \left( \omega \right) \right) }.
\end{equation}
Esta expresion puede simplificarse escribiendo
\begin{equation}
H\left( e^{-j \, \omega }\right) H\left( e^{j \, \omega }\right) =
\frac{1}{\left(1-a\cos \left( \omega \right) \right) ^{2}+
\left( a\sin \left( \omega\right) \right) ^{2}},
\end{equation}
y finalmente
\begin{equation}
H\left( e^{-j \, \omega }\right) H\left( e^{j \, \omega }\right) =
\frac{1}{1+a^{2}-2a\cos \left( \omega \right) }.
\end{equation}

\item  La fase de la respuesta en frecuencia es
\begin{equation}
\angle (H(e^{j \omega})) = -\arctan \left( \frac{a\sin \omega }{1-a\cos \left( \omega \right) }\right) .
\end{equation}
\end{enumerate}

\subsection{Sistema de tiempo discreto con dos polos}

Consideremos el sistema de tiempo discreto lineal, invariante en el tiempo y
causal cuya ecuaci\'{o}n de diferencias es
\begin{equation}
y\left[ n\right] -2r\cos \left( \theta \right) y\left[ n-1\right] +r^{2}y%
\left[ n-2\right] =x\left[ n\right] ,
\end{equation}
donde $0<r<1,$ y $0\leq \theta \leq \pi .$

\begin{enumerate}
\item  La funci\'{o}n de transferencia es
\begin{equation}
H\left( z\right) =\frac{1}{1-2r\cos \left( \theta \right) z^{-1}+r^{2}z^{-2}%
}=\frac{1}{\left( 1-re^{j\theta }z^{-1}\right) \left( 1-re^{-j\theta
}z^{-1}\right) }.
\end{equation}

\item  La respuesta en frecuencia es
\begin{equation}
H\left( e^{j \, \omega }\right) =\frac{1}{1-2r\cos \left( \theta \right)
e^{-j \, \omega }+r^{2}e^{-2j \, \omega }}=\frac{1}{\left( 1-re^{j\theta
}e^{-j \, \omega }\right) \left( 1-re^{-j\theta }e^{-j \, \omega }\right) }.
\end{equation}
\end{enumerate}

